\ifx \notincludehead\undefined
	\newcommand*{\No}{\textnumero}
	\documentclass[russian, utf8, 12pt, pointsubsection,floatsubsection]{eskdtext}
	\usepackage[russian]{babel}

	\ifx\pdfoutput\undefined  
		\def\pdfoutput{0}
	\fi

	\ifnum\pdfoutput=0 
		\sloppy
		%	\usepackage[dvips]{graphicx}       % загрузка графики под dvi  
		\textheight=250mm                  % для DVI высота печатного текста
		\textwidth=165mm                   % ширина печатного текста
	\else
		%	\usepackage[pdftex]{graphicx}      % загрузка графики под pdf
		\usepackage{cmap}                  % чтоб работал поиск по PDF 
		% гиперссылки в PDF
		\usepackage[unicode, pdftex, colorlinks, linkcolor=blue]{hyperref} 
		\pdfcompresslevel=9                % сжимаем PDF 
		%	\textheight=240mm                  % для PDF высота печатного текста
		%	\textwidth=165mm                   % ширина печатного текста
	\fi 

	\usepackage{eskdchngsheet}
	\usepackage[T2A]{fontenc}
	%\usepackage[cp1251]{inputenc}
	\usepackage{amstext}
	\usepackage{amsmath}
	\usepackage{listings}
	\usepackage{rotating}
	\usepackage{pmasc}

	\usepackage{placeins}    % пакет позволяет вставлять плавающие объекты (рисунки) в том месте, 
                         % где это необходимо. Для вывода рисунка после него встаить команду \FloatBarrier

    \usepackage{comment}
	\usepackage{array}
	\usepackage{longtable}   % подключение длинных таблиц
    \usepackage{multirow}    % для формирования таблиц с объединением строк
	\usepackage{indentfirst} % идентификация первых абзацев после секционирования
	\usepackage{fancyhdr}                    % расширенный формат страниц
	\usepackage{ulem}        % подчеркивания текста \uline\uuline\uwave\sout \xout
	%\voffset=-25mm   % -25                   % сдвиг страницы вверх
	%\hoffset=-15mm   % -10                   % сдвиг страницы влево

	\usepackage{floatflt} % для рисунков
	\usepackage{wrapfig}  % для рисунков
    \usepackage{xcolor}
    %\usepackage[most]{tcolorbox}
    \usepackage{subcaption} % Загружает пакет заголовков
    \usepackage{xspace}

	\sloppy                             % подавление дополнительных переносов
	\righthyphenmin=2                   % можно переносить
	
	\ESKDclassCode{ТП}
	\ESKDtitle{Программная система планирования производства ПС ПП «Opti-Corrugated» \FIRMA. Подключение гофроагрегата JSM}

	\usepackage{lscape}

	% для Первой спецификации
	\ESKDdocName{Технический рабочий проект\\ Пояснительная записка.}
	\ESKDsignature{65698922.425120.00036.П2.01.1.1}
	\ESKDcolumnII{65698922.425120.00036.П2.01.1.1}
	\ESKDcolumnI{Программная документация. ТЗ. Ред.1}

	\ESKDgroup{ООО <<Опти-Софт>>}
	\ESKDauthor{Косицын Д.П.}
	\ESKDtitleAgreedBy{Директор ООО <<Опти-Софт>>}{Шабаев А.И.}
	\ESKDtitleDesignedBy{Зам. директора <<Опти-Софт>>}{Косицын Д.П.}
    \ESKDtitleDesignedBy{Консультант}{Смирнов А.В.}	
    \ESKDtitleDesignedBy{Начальник отдела разработки гофротары}{Сошкин Р.В.}


	\ESKDtitleApprovedBy{Генеральный директор \FIRMA}{\DIRECTOR}
	%\ESKDtitleApprovedBy{\rule{72pt}{1pt}}{\rule{72pt}{1pt}}
	\ESKDtitleAgreedBy{Директор по производству}{ФИО}
	%\ESKDtitleAgreedBy{Коммерческий директор}{Фамилия И.О.}
% 	\ESKDtitleAgreedBy{Финансовый директор}{Ерофеева М.А.}
% 	\ESKDtitleAgreedBy{Начальник отдела автоматизации}{ФИО}

	%\ESKDtitleAgreedBy{\rule{72pt}{1pt}}{\rule{72pt}{1pt}}

	\ESKDdate{2026/03/21}

%	\newcommand*{\No}{\textnumero}

	% для нумерации в длинных enumerate: a, b,...y, z, aa,ab,..
	%\usepackage{alphalph}  
	%\renewcommand{\theenumi}{\alphalph{\value{enumi}}}
\fi



%\section {Назначение и цели создания системы}

\subsection{Назначение системы}
Проект глубокой модернизации программной системы планирования производства ПС ПП «Opti-Corrugated» для автоматизации процессов планирования и учета производства на \FIRMA. Подключение гофроагрегата JSM-3000.




\subsection{Цели и задачи работ по разработке ПС ПП}

В рамках данного проекта модернизации программного обеспечения преследуются следующие конкретные цели применительно к документу ''Выработка гофроагрегата'':

\begin{itemize}
    \item Автоматизация запуска производства. Исключить необходимость ручного ввода времени начала смены. При открытии формы машинист должен видеть только свои задания, привязанные к оборудованию его компьютера.
        \item Интеграция с гофроагрегатом. Убрать человеческий фактор по заведению заказов в память гофроагрегата.
    \item Повышение оперативности учета. Обеспечить автоматическую регистрацию выпуска продукции и полуфабрикатов на гофроагрегате и загрузку данных в систему \gofro в режиме реального времени.

    \item Повышение точности учёта. Перейти от минут к секундам во всех расчётах длительности простоя и скорости. Скорость рассчитывать строго по чистому времени работы оборудования (исключив время пуско-наладочных работ).
    \item Обеспечение целостности данных. Добавить проверки заполнения обязательных полей перед сохранением или завершением смены. Блокировать сохранение, если не указана причина простоя или брака.
    \item Унификация интерфейсов. Реализовать единый механизм выбора причины возникновения простоя или выпуска брака с обязательным указанием пояснения. % Для быстрого добавления единичных случаев брака разработать отдельную форму.
    \item Оптимизация командной панели. На основной форме вывести кнопки для старта и завершения этапов работы (переналадка, простой), чтобы машинист мог управлять процессом вручную при необходимости.

\end{itemize}

Эти цели направлены на снижение ошибок ввода, повышение прозрачности производственного процесса и обеспечение автоматического контроля качества данных.

